% Bagian: first-order-logic
% Bab: models-theories
% Subbagian: expressing-relations

\documentclass[../../../include/open-logic-section]{subfiles}

\begin{document}

\olfileid[id]{fol}{mat}{exr}

\olsection{Mengungkapkan Relasi dalam \article{structure}
  \printtoken{S}{structure}}

\begin{explain}
Salah satu kegunaan utama !!{formula}s ialah mengungkapkan sifat dan relasi
dalam !!a{structure}~$\Struct M$ berdasarkan konsep-konsep primitif bahasa
$\Lang L$ dari~$\Struct M$. Maksudnya adalah sebagai berikut: !!{domain} dari
$\Struct M$ merupakan suatu himpunan objek. !!^{constant}s, !!{function}s,
dan !!{predicate}s ditafsirkan dalam~$\Struct M$ masing-masing oleh beberapa
objek dalam~$\Domain M$, fungsi pada~$\Domain M$, dan relasi pada~$\Domain M$.
Misalnya, jika $\Obj A^2_0$ berada dalam $\Lang L$, maka $\Struct M$
menetapkan kepadanya suatu relasi~$R = \Assign{{\Obj A^2_0}}{M}$. Kemudian
formula $\Atom{\Obj A^2_0}{\Obj v_1, \Obj v_2}$ \emph{mengungkapkan} relasi
itu sendiri dalam pengertian berikut: jika suatu penugasan variabel~$s$
memetakan $\Obj v_1$ ke $a \in \Domain{M}$ dan $\Obj v_2$ ke
$b \in \Domain M$, maka
\[
Rab \text{\quad jika dan hanya jika\quad} \Sat{M}{\Atom{\Obj A^2_0}{\Obj v_1, \Obj v_2}}[s].
\]
Perhatikan bahwa kita harus melibatkan penugasan variabel di sini: kita tidak
dapat sekadar mengatakan ``$Rab$ jika dan hanya jika
$\Sat{M}{\Atom{\Obj A^2_0}{a, b}}$'' karena $a$ dan $b$ bukan simbol bahasa
kita; keduanya merupakan !!{element}s dari~$\Domain{M}$.

Karena kita tidak hanya mempunyai !!{formula}s atomik, tetapi juga dapat
menggabungkannya dengan penghubung logis dan kuantor, !!{formula}s yang lebih
rumit dapat mendefinisikan relasi lain yang tidak dibangun langsung ke
dalam~$\Struct M$. Kita tertarik pada cara melakukan hal tersebut dan,
khususnya, pada relasi mana yang dapat kita definisikan dalam !!a{structure}.
\end{explain}

\begin{defn}
Misalkan $!A(\Obj v_1,\dots, \Obj v_n)$ suatu !!a{formula} dari $\Lang L$
yang hanya mempunyai $\Obj v_1$,\dots, $\Obj v_n$ sebagai variabel bebas,
dan misalkan $\Struct M$ suatu !!a{structure} bagi~$\Lang L$.
$!A(\Obj v_1,\dots, \Obj v_n)$ \emph{mengungkapkan relasi}~$R \subseteq
\Domain M^n$ jika dan hanya jika
\[
Ra_1\dots a_n \text{\quad jika dan hanya jika\quad} \Sat{M}{\Atom{!A}{\Obj
    v_1,\dots, \Obj v_n}}[s]
\]
bagi sebarang penugasan variabel~$s$ dengan $s(\Obj v_i) = a_i$ ($i = 1,
\dots, n$).
\end{defn}

\begin{ex}
Dalam model standar aritmetika~$\Struct N$, !!{formula} $\Obj v_1 < \Obj v_2
\lor \eq[\Obj v_1][\Obj v_2]$ mengungkapkan relasi $\le$ pada~$\Nat$.
!!^{formula} $\eq[\Obj v_2][\Obj v_1']$ mengungkapkan relasi penerus, yakni
relasi $R \subseteq \Nat^2$ dengan $Rnm$ berlaku jika $m$ merupakan penerus
dari~$n$. Formula $\eq[\Obj v_1][\Obj v_2']$ mengungkapkan relasi pendahulu.
!!^{formula}s $\lexists[\Obj v_3][(\eq/[\Obj v_3][\Obj 0] \land
  \eq[\Obj v_2][(\Obj v_1 + \Obj v_3)])]$ dan $\lexists[\Obj
  v_3][\eq[(\Obj v_1 + {\Obj v_3}')][\Obj v_2]]$ keduanya mengungkapkan relasi
$\Obj <$. Ini berarti bahwa simbol predikat~$<$ sebenarnya berlebihan dalam
bahasa aritmetika; simbol itu dapat didefinisikan.
\end{ex}

\begin{explain}
Gagasan ini tidak hanya menarik dalam !!{structure}s tertentu, tetapi juga
secara umum setiap kali kita menggunakan suatu bahasa untuk menggambarkan satu
atau beberapa model yang dimaksudkan, yakni ketika kita meninjau teori. Teori
semacam itu sering kali hanya memuat sedikit !!{predicate}s sebagai simbol
dasar, tetapi banyak relasi lain sering berperan penting dalam domain yang
hendak dideskripsikan. Jika relasi-relasi lain tersebut dapat diungkapkan
secara sistematis oleh relasi-relasi yang menafsirkan !!{predicate}s dasar
bahasa itu, kita mengatakan bahwa relasi-relasi tersebut dapat
\emph{didefinisikan} dalam bahasa tersebut.
\end{explain}

\begin{prob}
Carilah !!{formula}s dalam $\Lang L_A$ yang mendefinisikan relasi-relasi
berikut:
\begin{enumerate}
\item $n$ berada di antara $i$ dan $j$;
\item $n$ membagi habis $m$ (yakni, $m$ merupakan kelipatan $n$);
\item $n$ merupakan bilangan prima (yakni, tidak ada bilangan selain $1$ dan
  $n$ yang membagi habis~$n$).
\end{enumerate}
\end{prob}

\begin{prob}
Andaikan formula $!A(\Obj v_1, \Obj v_2)$ mengungkapkan relasi
$R \subseteq \Domain M^2$ dalam !!a{structure}~$\Struct M$. Carilah formula
yang mengungkapkan relasi-relasi berikut:
\begin{enumerate}
\item invers $R^{-1}$ dari $R$;
\item hasil kali relatif $R \mid R$;
\end{enumerate}
Dapatkah Anda menemukan cara untuk mengungkapkan $R^+$, penutupan transitif
dari~$R$?
\end{prob}

\begin{prob}
Misalkan $\Lang{L}$ bahasa yang hanya memuat simbol predikat dua-tempat $<$
(tanpa !!{constant}s, !!{function}s, atau !!{predicate}s lain---kecuali,
tentu saja,~$\eq$). Misalkan $\Struct{N}$ struktur sedemikian sehingga
$\Domain{N} = \Nat$ dan $\Assign{<}{N} = \Setabs{\tuple{n,m}}{n < m}$.
Buktikan hal-hal berikut:
\begin{enumerate}
\item $\{ 0 \}$ dapat didefinisikan dalam $\Struct{N}$;
\item $\{ 1 \}$ dapat didefinisikan dalam $\Struct{N}$;
\item $\{ 2 \}$ dapat didefinisikan dalam $\Struct{N}$;
\item bagi setiap $n \in \Nat$, himpunan $\{ n \}$ dapat didefinisikan
  dalam $\Struct{N}$;
\item setiap himpunan bagian hingga dari $\Domain{N}$ dapat didefinisikan
  dalam $\Struct{N}$;
\item setiap himpunan bagian kohingga dari $\Domain{N}$ dapat didefinisikan
  dalam $\Struct{N}$ (dengan $X \subseteq \Nat$ disebut kohingga jika dan
  hanya jika $\Nat \setminus X$ hingga).
\end{enumerate}
\end{prob}


\end{document}
